\section{过滤：粗糙但高吞吐的打分器}

\subsection{问题的本质}

把过滤写成一个公式化问题并不复杂。给定目标数据 $T$ 和原始数据 $R$，我们想要一个子集 $R' \subset R$，让它在某个意义下更接近 $T$。关键不在于“精确拟合”，而在于：

\begin{itemize}
    \item 你要学会对样本打分；
    \item 你要把高分样本留住；
    \item 你要在超大规模数据上完成这个动作。
\end{itemize}

\begin{warningbox}{为什么不能直接上大模型}
如果过滤器本身和训练模型差不多重，那么你在筛掉 99\% 数据之前，就已经为这 99\% 数据付出了接近训练级别的算力。那样做通常不划算。
\end{warningbox}

\begin{figure}[H]
\centering
\begin{tikzpicture}[node distance=1.1cm, every node/.style={font=\small, align=center}]
\node[draw, rounded corners, fill=red!8, minimum width=2.7cm, minimum height=0.8cm] (bad) {低分样本};
\node[draw, rounded corners, fill=yellow!14, right=1.7cm of bad, minimum width=2.7cm, minimum height=0.8cm] (mid) {边界样本};
\node[draw, rounded corners, fill=green!14, right=1.7cm of mid, minimum width=2.7cm, minimum height=0.8cm] (good) {高分样本};
\draw[-{Latex[length=3mm]}, thick] (bad) -- (mid);
\draw[-{Latex[length=3mm]}, thick] (mid) -- (good);
\node[below=0.8cm of mid, align=center] {阈值越高，保留越少\\但质量通常越高};
\end{tikzpicture}
\caption{过滤的核心动作是阈值化或重采样：把连续分数转成是否保留的决定}
\end{figure}

\subsection{通用过滤模板}

\begin{lstlisting}[language=Python, caption={一个极简的过滤模板}]
model = fit(target_data)
scores = [model.score(x) for x in raw_data]
keep = [x for x, s in zip(raw_data, scores) if s >= threshold]
\end{lstlisting}

这个模板几乎就是本讲所有方法的共同骨架。

差异只在于打分函数 \texttt{score(x)} 是什么，以及 \texttt{fit(target\_data)} 用什么模型来实现。

\subsection{分数、阈值和重采样}

同一个分数，在工程里可以有三种常见用法。

\begin{table}[H]
\centering
\small
\renewcommand{\arraystretch}{1.15}
\begin{tabular}{p{2.2cm}p{4.3cm}p{4.1cm}p{3.1cm}}
\toprule
\textbf{用法} & \textbf{动作} & \textbf{适合场景} & \textbf{代价} \\
\midrule
threshold & 高于阈值就保留 & 噪声过滤、毒性过滤 & 容易受校准影响 \\
ranking & 取 top-$k$ 或 top-$p$\% & CCNet 风格排序、预算固定 & 结果依赖样本总量 \\
resampling & 按权重重采样 & DSIR、分布匹配 & 统计波动更大 \\
\bottomrule
\end{tabular}
\caption{同一个分数函数，可以驱动三种不同的数据选择动作}
\end{table}

\begin{importantbox}{为什么要区分这三种动作}
阈值适合“要还是不要”的粗过滤；排序适合固定预算下的高质量截断；重采样适合想让样本整体分布更像目标分布的场景。把它们混在一起，往往会导致目标不清晰。
\end{importantbox}

\begin{figure}[H]
\centering
\begin{tikzpicture}[node distance=1.0cm, every node/.style={font=\small, align=center}]
\node[draw, rounded corners, fill=blue!10, minimum width=2.8cm, minimum height=0.8cm] (score) {score all samples};
\node[draw, rounded corners, fill=yellow!16, below left=1.0cm and 1.2cm of score, minimum width=2.8cm, minimum height=0.8cm] (thr) {threshold};
\node[draw, rounded corners, fill=orange!14, below=1.0cm of score, minimum width=2.8cm, minimum height=0.8cm] (rank) {rank and keep top-$k$};
\node[draw, rounded corners, fill=green!14, below right=1.0cm and 1.2cm of score, minimum width=2.8cm, minimum height=0.8cm] (res) {weighted resample};
\draw[-{Latex[length=3mm]}, thick] (score) -- (thr);
\draw[-{Latex[length=3mm]}, thick] (score) -- (rank);
\draw[-{Latex[length=3mm]}, thick] (score) -- (res);
\end{tikzpicture}
\caption{同一组分数可以派生出三种选择策略：阈值、排序和重采样}
\end{figure}

